\PassOptionsToPackage{table, dvipsnames}{xcolor} %In pratica dice a LaTeX: “Quando carichi il pacchetto xcolor, passa queste opzioni”. table permette di colorare anche le tabelle (es. righe o colonne). dvipsnames abilita un set di ~68 colori predefiniti con nomi leggibili.xcolor è un pacchetto LaTeX per gestire i colori in documenti, figure, tabelle, grafici, testi, ecc
\documentclass[a4paper,twoside,12pt]{toptesi}
\usepackage[T1]{fontenc}
\usepackage[utf8]{inputenc}
\usepackage[italian]{babel} %Gestione e codifica della lingua italiana
\usepackage{graphicx}  % Required for inserting images
\usepackage{tikz} %per flowchart, diagrammi di flusso o modelli di process
\usetikzlibrary{positioning}  %Senza positioning dovresti calcolare coordinate (x,y) manualmente, il che è lungo e soggetto a errori
\usepackage{pgfplots}  %PGFPlots è costruito sopra TikZ, specifico per grafici di dati e funzioni matematiche.
\usepackage[newfloat]{minted}  %permette di inserire codice sorgente evidenziato (syntax highlighting) all’interno del documento. Supporta tantissimi linguaggi: Python, Java, C++, SQL, HTML, ecc. [newfloat permette di trattare i blocchi di codice come veri e propri float, cioè oggetti simili a figure o tabelle. \SetupFloatingEnvironment{listing}{name=Codice}]
\usepackage[font=small,skip=1pt]{caption}  %Serve per personalizzare l’aspetto delle didascalie di figure, tabelle e, se usi minted[newfloat], anche dei blocchi di codice (float)
\captionsetup[listing]{position=top} %Dice a LaTeX di mettere la didascalia sopra il blocco di codice, invece che sotto (default).
\usepackage{subcaption} %subcaption è un pacchetto LaTeX che permette di gestire sottodidascalie all’interno di figure o tabelle. Serve quando vuoi avere più immagini o grafici in un unico blocco, ognuno con la sua didascalia, ma comunque raggruppati sotto un’unica figura principale.
\pgfplotsset{compat=1.18}

\SetupFloatingEnvironment{listing}{name=Codice} %“Codice 1, Codice 2…” è più leggibile e professionale per una tesi italiana rispetto a “Listing 1…”

\newcommand{\walter}[1]{\textcolor{red}{{\bf [Walter: }{#1}{\bf ]}}}

\usepackage{hyperref} %Aggiunge link cliccabili

%questo rende cliccabili i link, l’indice, le citazioni, ecc.
\hypersetup{
    colorlinks = true,
    allcolors=black
}

%%% workaround for listing
% save the current meaning of \listing
\let\savedlisting\listing
% at the right spot, restore the meaning
\AtBeginDocument{\let\listing\savedlisting}
%%% end of workaround


\def\dept{DIPARTIMENTO DI INGEGNERIA ELETTRICA E DELL'INFORMAZIONE}
\def\course{Corso di Laurea in Ingegneria Informatica e dell’Automazione}
\def\title{Titolo della tesi}
\def\author{Nicholas Agrusti}
\def\prof{Prof. Ing. Ciccio Cappuccio}
\def\subject{Intelligenza Artificiale}
\def\annoacc{2022 - 2023}
\def\beforecandidate{Laureando}
\def\beforetitle{Tesi di Laurea in}
\def\beforeprof{Relatore}
\def\beforeannoacc{Anno Accademico}
%Queste variabili servono per generare automaticamente il frontespizio

\makeatletter
\def\cleardoublepage{\clearpage\if@twoside \ifodd\c@page\else
    \hbox{}
    \vspace*{\fill}
    \vspace{\fill}
    \thispagestyle{empty}
    \newpage
    \if@twocolumn\hbox{}\newpage\fi\fi\fi}
\makeatother
%In LaTeX, il comando \cleardoublepage serve a terminare la pagina corrente e, se il documento è a due facciate (twoside), fa in modo che la pagina successiva inizi sempre su destra (pagina dispari


\begin{document}

%\maketitle

\begingroup
\thispagestyle{empty}
\begin{titlepage}
	\begin{tikzpicture}[remember picture,overlay]
		\centering
			\node[yshift=-6 cm] (logo) at (current page.north) {\includegraphics[width=0.35\linewidth]{images/poliba.png}};
			\node[text width=50em,yshift=0.25cm, align = center, below = of logo](poliba){\bfseries \Large POLITECNICO DI BARI};
			\node[text width=40em, align = center, yshift=.55cm,below = of poliba](course){\normalsize \dept \\
%
				\normalsize \textbf{\course}};
		\node[text width=35em,align = center,  yshift=1.2cm,below = of course](line){\par\noindent\rule{\textwidth}{0.4pt}};
		\node[text width=40em, align = center, yshift=.55cm,below = of line](lia){\beforetitle\xspace \subject};
		\node[text width=40em, align = center, yshift=-0.5cm,below = of lia](title){\bfseries \fontsize{21pt}{20pt}\selectfont \title\par};
	\node[text width=35em, align = left, yshift=-0.75cm,below = of title](supervisor){\large \beforeprof \\ \textbf{\prof}};
		\node[text width=35em, align = right, yshift=0cm,below = of supervisor](candidate){\large \beforecandidate\\ \textbf{\author}};
		\node[text width=35em,align = center,  yshift= 0cm,below = of candidate](line2){\par\noindent\rule{\textwidth}{0.4pt}};
		\node[text width=50em, align = center, yshift=0.5cm,below = of line2](year){\beforeannoacc\xspace \annoacc};
	\end{tikzpicture}
\end{titlepage}       % frontespizio
\endgroup

\cleardoublepage
\input{dedica} %%include sezioni piccole (senza cambiare pagina)

\cleardoublepage

\pagenumbering{roman}

\input{abstract} %Presentare l’argomento principale della tesi. Riassumere obiettivi, metodo e risultati principali.

\cleardoublepage

\tableofcontents %genera automaticamente l’indice

\cleardoublepage

\pagenumbering{arabic}
\setcounter{page}{1}
%\include{introduzione}
%\include{capitolo_1}
%\include{capitolo_2}
%\include{capitolo_3}
%\include{capitolo_4}
%\include{conclusione}
%\include{ringraziamenti}


\bibliographystyle{plain}
\bibliography{bibliografia} %crea la bibliografia finale.
\end{document}